% =====================================================================
%  1bat-ud6-2.tex  —  SESSIÓ 2: La regla de Laplace i les propietats bàsiques
%  (sense preàmbul: s'inclou des de main.tex)
% =====================================================================
\horatitol{Sessió 2 — La regla de Laplace i les propietats bàsiques}%
{Posem nom i fórmula a la intuïció d'ahir: la regla de Laplace, les propietats de la probabilitat i el succés contrari.}

\subsection{Idea clau}

A la sessió 1 vam intuir que la probabilitat és «casos favorables sobre casos
possibles». Avui ho \textbf{formalitzem} amb la \textbf{regla de Laplace} i fixem les
\textbf{propietats bàsiques} que farem servir tota la unitat, treballant sempre les
tres formes (fracció, decimal i percentatge) que demana el criteri d'avaluació.

\begin{idea}
\textbf{Regla de Laplace i propietats}\\
Si tots els resultats són \textbf{equiprobables} (tenen la mateixa possibilitat):
\[
P(A)=\frac{\text{nombre de casos favorables a }A}{\text{nombre de casos possibles}}.
\]
\textbf{Propietats bàsiques:}
\begin{itemize}[leftmargin=1.4em, itemsep=2pt]
  \item $0\le P(A)\le 1$ sempre (cap probabilitat surt d'aquest interval).
  \item $P(\text{succés segur})=1$ \quad i \quad $P(\text{succés impossible})=0$.
  \item \textbf{Succés contrari} ($\overline{A}$, «que \emph{no} passi $A$»):
        \[
        P(\overline{A})=1-P(A).
        \]
\end{itemize}
\textbf{Recorda:} $\frac{1}{4}=0,25=25\%$ són la mateixa probabilitat.
\end{idea}

\subsection{Exemples}

\begin{exemple}[Laplace amb una baralla]
En una baralla espanyola de $40$ cartes, quina és la probabilitat de treure un rei?
Casos possibles: $40$. Casos favorables: $4$ (hi ha $4$ reis). Per la regla de Laplace:
\[
P(\text{rei})=\frac{4}{40}=\frac{1}{10}=0,1=10\%.
\]
I la probabilitat de \textbf{no} treure rei (succés contrari)?
\[
P(\overline{\text{rei}})=1-0,1=0,9=90\%.
\]
\end{exemple}

\begin{exemple}[el succés contrari estalvia feina]
En tirar dos daus, calcular la probabilitat de treure \emph{almenys} un $6$ comptant
tots els casos és llarg. És molt més curt el \textbf{contrari}: «\emph{cap} $6$». La
probabilitat de no treure $6$ en un dau és $\frac{5}{6}$; en dos daus,
$\frac{5}{6}\per\frac{5}{6}=\frac{25}{36}$. Per tant:
\[
P(\text{almenys un 6})=1-\frac{25}{36}=\frac{11}{36}\approx 0,31=31\%.
\]
Quan un enunciat diu «almenys», sovint convé passar pel succés contrari.
\end{exemple}

\subsection{Exercicis resolts}

\begin{exresolt}[probabilitat i contrari sobre un grup d'usuaris]
En un casal hi ha $25$ usuaris: $15$ són dones i $10$, homes. Si en triem un a l'atzar,
quina és la probabilitat que sigui dona? I que sigui home? Expressa-les en les tres
formes.
\solucio
\textbf{Dona:} casos favorables $15$, possibles $25$:
\[
P(\text{dona})=\frac{15}{25}=\frac{3}{5}=0,6=60\%.
\]
\textbf{Home} (succés contrari de dona):
\[
P(\text{home})=1-0,6=0,4=40\% \quad\Big(\text{o bé }\tfrac{10}{25}=\tfrac{2}{5}\Big).
\]
\resposta{$P(\text{dona})=\frac{3}{5}=0,6=60\%$; $P(\text{home})=\frac{2}{5}=0,4=40\%$.}
\end{exresolt}

\begin{exresolt}[comprovar amb les propietats]
En una bossa hi ha boles numerades de l'$1$ al $20$. Calcula la probabilitat de treure
un múltiple de $5$ i comprova que el resultat és coherent amb les propietats.
\solucio
Múltiples de $5$ entre $1$ i $20$: el $5$, el $10$, el $15$ i el $20$ ($4$ casos
favorables). Casos possibles: $20$.
\[
P(\text{múltiple de 5})=\frac{4}{20}=\frac{1}{5}=0,2=20\%.
\]
\textbf{Comprovació:} $0,2$ està entre $0$ i $1$ (correcte). El contrari («no múltiple de
$5$») és $1-0,2=0,8=80\%$, que també és un valor vàlid.
\resposta{$\frac{1}{5}=0,2=20\%$; coherent (entre $0$ i $1$).}
\end{exresolt}

\subsection{Exercicis per fer}

\begin{exercicis}
\begin{exlist}
  \item En una baralla de $40$ cartes, calcula (en fracció, decimal i percentatge) la
        probabilitat de treure: (a) una copa (n'hi ha $10$); (b) un as (n'hi ha $4$);
        (c) una carta que \textbf{no} sigui copa.
  \item En tirar un dau, calcula la probabilitat de cada succés i la del seu contrari:
        (a) treure un $5$; (b) treure un nombre parell.
  \item En una bossa amb boles numerades de l'$1$ al $30$, calcula la probabilitat de
        treure un múltiple de $3$. Expressa-la en les tres formes.
  \item En un grup de $50$ persones, $35$ tenen el carnet de la biblioteca. Quina és la
        probabilitat que una persona triada a l'atzar el tingui? I que no el tingui?
  \item \textbf{(Almenys)} En tirar dues monedes, calcula la probabilitat de treure
        \textbf{almenys} una cara (pista: passa pel contrari, «cap cara»).
  \item \textbf{(Raonament)} Per què la regla de Laplace només val quan els casos són
        \textbf{equiprobables}? Posa un exemple de situació on \emph{no} ho siguin i, per
        tant, no es pugui aplicar directament.
\end{exlist}
\end{exercicis}

% ---------------------------------------------------------------------
\fullsolucions{Sessió 2}
\begin{solucions}
\begin{sollist}
  \item (a) $\frac{10}{40}=\frac{1}{4}=0,25=25\%$; (b) $\frac{4}{40}=\frac{1}{10}=0,1=10\%$;
        (c) no copa: $1-0,25=0,75=75\%$.
  \item (a) $P(5)=\frac{1}{6}\approx 0,17$; contrari $\frac{5}{6}\approx 0,83$. \quad
        (b) $P(\text{parell})=\frac{3}{6}=\frac{1}{2}=0,5$; contrari $0,5$.
  \item Múltiples de $3$ fins a $30$: $3,6,9,\dots,30$, en total $10$.
        $P=\frac{10}{30}=\frac{1}{3}\approx 0,33=33,3\%$.
  \item $P(\text{té carnet})=\frac{35}{50}=\frac{7}{10}=0,7=70\%$; no el té:
        $1-0,7=0,3=30\%$.
  \item Contrari «cap cara» (les dues creus): $\frac{1}{2}\per\frac{1}{2}=\frac{1}{4}$.
        Per tant $P(\text{almenys una cara})=1-\frac{1}{4}=\frac{3}{4}=0,75=75\%$.
  \item Perquè Laplace reparteix la probabilitat \textbf{a parts iguals} entre tots els
        casos; si uns són més probables que altres, dividir favorables entre possibles
        dóna un resultat fals. Exemple: una xinxeta que cau de punta o de cap \emph{no}
        té $\frac{1}{2}$ de cada (una posició és més probable); o una ruleta trucada.
\end{sollist}
\end{solucions}
